\documentclass[12pt]{article} 

\usepackage[top=1.5cm, bottom=1.5cm, left=1.5cm, right=1.5cm]{geometry}
\usepackage{color,graphicx,amsmath,mathtools,amssymb}
\pagestyle{plain}

\begin{document}  
\pagestyle{empty}
 



%
%\vskip0.5in
%
%
%\noindent {\small This document is property of Northeastern University.  Unauthorized distribution of the content is not~permitted.}

Homework 3: Nico Boyle, Owen Moorhus, Asher Olgeirson

\begin{enumerate}


% \item

%       %%%%   Omitting the problem statement to reduce clutter      %%%%

% A league has 16 teams.   In the regular season each team plays each other 6 times.
% The season is played in 90 ``rounds". In  each round every team plays a game. 
% Each team's  rank is determined by number of wins.   If two teams end up with the same record, their rank is determined via a new tie-breaker game. (What's relevant is that teams that end up with the same record at the end of the regular season have unresolved relative ranks).

% During the season every team does their best to increase their rank, until possibly reaching a time when they realize that their rank can't change.

% $\bullet$ What is the earliest possible round at which all teams in the league no longer care about improving their rank?  Equivalently, how large could $r$ be, so that with $r$ regular season rounds remaining none of the teams care? \\
% $\bullet$ You should include an example of all team win records and justify why all teams no longer care.  \\
% $\bullet$  You should explain why your value of $r$ couldn't be made larger.  Ideally you should mention the pigeonhole principle.\\


\item

{\color{blue}
Largest valid value of $r$:  \\

$r=4$
\vskip1cm
Example of win records: \\

81,76,71,66,61,56,51,46,40,35,30,25,20,15,10,5\\

Justification: \\

The teams must be ranked s.t. $W_1 > W_2 > W_3 > ... > W_{16}$, where $W_i$ is the number of wins of team $i$.  \\
16 total teams means there are 15 gaps between the teams.\\
For no team to care about improving their rank, the gap between each team must be at least r+1 which leads to: 
$W_1 - W_{16} \geq 15(r+1)$\\
After 90 - r rounds, the largest possible difference between the first and last team is: 
$90-r \geq 15(r+1)$\\
Solving this inequality for r gives $r \leq 4.6875$.\\
And because r must be an integer, $r=4$ is the largest valid value of r.\\

\vskip1cm
The pigeonhole principle is shows why r cannot be made larger than 4.\\
If r were 5, with 15 gaps between the teams, there would need to be a total spread of 90: $15\cdot 6 = 90$. \\
But with only 5 rounds remaining, the maximum spread is 85.\\
Thus, if we take the possible win totals of the teams as the holes and the teams as the pigeons, then, with 5 rounds remaining,
there are only 86 possible win totals (0 through 85).\\
For no team to care about improving their rank, team wins must differ by at least 6, which means that there can only be 15 possible win totals (0,6,12,...,90).\\
Thus, there are not enough holes to place all 16 teams while maintaing the required spacing of wins.\\
}







%%%%%%%%%%%%%%%%%%%%%%%%%%%%%%%%%%%%%%%%%
\newpage

\item
Simplify the following expression, making one modification per line.\\
On each line you should also  briefly explain what rule was used to get that line from the one above. 
$\neg ( \neg (p\land q) \lor \neg (\neg q \land r) )$ \\

%%%%.  To avoid LaTeX errors, make sure to leave a space after \neg \land \lor as in the sample. (E.g., don't write \negq)

{\color{blue}
1. $\neg ( (\neg p \lor \neg q) \lor \neg (\neg q \land r) ) \rightarrow$ Demorgan's on $\neg (p\land q) $\\
2. $\neg ( (\neg p \lor \neg q) \lor  (q \lor \neg r) ) \rightarrow$ Demorgan's on $ \neg (\neg q \land r) $\\
3. $\neg ( \neg p \lor \neg q \lor  q \lor \neg r ) \rightarrow$ Associative Law \\
4. $p \land q \land  \neg q \land r \rightarrow $ Demorgan's on outer parenthesis \\
5. $p \land F \land r \rightarrow$ $q \land  \neg q$ annihilates to F \\
6. $F \rightarrow$ $p \land F \land r$ is always F \\
}

%%%%%%%%%%%%%%%%%%%%%%%%%%%%%%%%%%%%%%%%%
\newpage


\item       %%%%.     Removing part of the problem statement to reduce clutter     %%%%%
%
%
% Let $X$ denote a logic operator, defined as 
Definition:    $a X b = \neg ( a \lor  b)$\\\\
Show that we can simulate the following three expressions, with other expressions in which the only operator used is $X$. The expressions shouldn't include T, F or negation either.  You should provide step-by-step derivation with comments. Each change should be on a new line. Each line should be equivalent to the preceding line. \\
a) $\neg p$ \\

{\color{blue}

$\neg p \leftrightarrow \neg p \lor \neg p$ as $p\lor p$ is always p

$\neg p \lor \neg p \leftrightarrow \neg (p\lor p)$ by deMorgan's laws

$\neg(p\lor p)\leftrightarrow pXp$ by definition

}

\vskip1cm

b) $p\lor q$ \\

{\color{blue}
  $p\lor q \iff \neg (\neg (p\lor q))$ a double logical not cancels out via deMorgan's laws\\ 
  $\iff \neg(pXq)$ the X operator is defined with a logical not \\ 
  $\iff (pXq)X(pXq)$ as $pXp \iff \neg p$ as proved above

}

\vskip1cm

c) $p\land q$\\

{\color{blue}
$p\land q \leftrightarrow \neg((\neg(p\lor p))\lor(\neg(q\lor q))$ assumed for the sake of argument \\
$(\neg p) X (\neg q) \leftrightarrow \neg (\neg p \lor \neg q) $ as $pXp \leftrightarrow \neg p$, which was proven above \\
$\neg(\neg p \lor \neg q) \leftrightarrow p\land q $ via truth table as:\\\\

\[
\begin{array}{|c|c|c|}
\hline
p & q & \neg(\neg p \lor \neg q) \\
\hline
T & T & T \\
T & F & F \\
F & T & F \\
F & F & F \\
\hline
\end{array}
\quad \leftrightarrow \quad
\begin{array}{|c|c|c|}
\hline
p & q & p \land q \\
\hline
T & T & T \\
T & F & F \\
F & T & F \\
F & F & F \\
\hline
\end{array}
\]

}



\end{enumerate} 




\end{document}
